% Símbolos Unicode no corpo (pdflatex) — português via utf8/inputenc
\ifxetex
\else
  \usepackage{newunicodechar}
  \newunicodechar{≥}{$\geq$}
  \newunicodechar{≤}{$\leq$}
  \newunicodechar{≠}{$\neq$}
  \newunicodechar{≈}{$\approx$}
  \newunicodechar{→}{$\rightarrow$}
  \newunicodechar{←}{$\leftarrow$}
  \newunicodechar{↔}{$\leftrightarrow$}
  \newunicodechar{−}{$-$}
  \newunicodechar{—}{---}
  \newunicodechar{–}{--}
  \newunicodechar{“}{``}
  \newunicodechar{”}{''}
  \newunicodechar{‘}{`}
  \newunicodechar{’}{'}
  \newunicodechar{…}{\ldots}
  \newunicodechar{″}{''}
  \newunicodechar{×}{\times}
  \newunicodechar{©}{\textcopyright}
  \newunicodechar{§}{\S}
  \newunicodechar{ª}{\textordfeminine}
  \newunicodechar{º}{\textordmasculine}
  \newunicodechar{┌}{}
  \newunicodechar{└}{}
  \newunicodechar{│}{|}
  \newunicodechar{─}{-}
  \newunicodechar{▲}{}
  \newunicodechar{▼}{}
  \newunicodechar{►}{}
\fi
